%%%%%%%%%%%%%%%%%%%%%%%%%%%%%%%%%%%%%%%%%
% Beamer Presentation
% LaTeX Template
% Version 1.0 (10/11/12)
%
% This template has been downloaded from:
% http://www.LaTeXTemplates.com
%
% License:
% CC BY-NC-SA 3.0 (http://creativecommons.org/licenses/by-nc-sa/3.0/)
%
%%%%%%%%%%%%%%%%%%%%%%%%%%%%%%%%%%%%%%%%%

%----------------------------------------------------------------------------------------
%	PACKAGES AND THEMES
%----------------------------------------------------------------------------------------

\documentclass{beamer}

\mode<presentation> {

% The Beamer class comes with a number of default slide themes
% which change the colors and layouts of slides. Below this is a list
% of all the themes, uncomment each in turn to see what they look like.

%\usetheme{default}
%\usetheme{AnnArbor}
%\usetheme{Antibes}
%\usetheme{Bergen}
%\usetheme{Berkeley}
%\usetheme{Berlin}
%\usetheme{Boadilla}
%\usetheme{CambridgeUS}
%\usetheme{Copenhagen}
%\usetheme{Darmstadt}
%\usetheme{Dresden}
%\usetheme{Frankfurt}
%\usetheme{Goettingen}
%\usetheme{Hannover}
%\usetheme{Ilmenau}
%\usetheme{JuanLesPins}
%\usetheme{Luebeck}
%\usetheme{Madrid}
%\usetheme{Malmoe}
%\usetheme{Marburg}
%\usetheme{Montpellier}
%\usetheme{PaloAlto}
\usetheme{Pittsburgh}
%\usetheme{Rochester}
%\usetheme{Singapore}
%\usetheme{Szeged}
%\usetheme{Warsaw}

% As well as themes, the Beamer class has a number of color themes
% for any slide theme. Uncomment each of these in turn to see how it
% changes the colors of your current slide theme.

%\usecolortheme{albatross}
\usecolortheme{beaver}
%\usecolortheme{beetle}
%\usecolortheme{crane}
%\usecolortheme{dolphin}
%\usecolortheme{dove}
%\usecolortheme{fly}
%\usecolortheme{lily}
%\usecolortheme{orchid}
%\usecolortheme{rose}
%\usecolortheme{seagull}
%\usecolortheme{seahorse}
%\usecolortheme{whale}
%\usecolortheme{wolverine}

%\setbeamertemplate{footline} % To remove the footer line in all slides uncomment this line
%\setbeamertemplate{footline}[page number] % To replace the footer line in all slides with a simple slide count uncomment this line

%\setbeamertemplate{navigation symbols}{} % To remove the navigation symbols from the bottom of all slides uncomment this line
}
\usepackage{graphicx} % Allows including images
\usepackage{booktabs} % Allows the use of \toprule, \midrule and \bottomrule in tables

\usepackage{fontspec}
\setsansfont{Calibri}[
    Path=Font/,
    Scale=0.9,
    Extension = .ttf,
    UprightFont=*-Regular,
    BoldFont=*-Bold,
    ItalicFont=*-Italic,
    BoldItalicFont=*-BoldItalic
    ]
\usepackage{unicode-math}
\setmathfont{Libertinus Math}% 使用默认的数学字体或你偏好的数学字体 Latin Modern Math/Libertinus Math
% \setmathfont[version=bold]{Latin Modern Math}

\usepackage{ragged2e}
\justifying\let\raggedright\justifying %两端对齐
%----------------------------------------------------------------------------------------
%	TITLE PAGE
%----------------------------------------------------------------------------------------

\title[Asset Pricing]{\textbf{Asset Pricing}} % The short title appears at the bottom of every slide, the full title is only on the title page

\subtitle{7. Factor Zoo}

\author[Cheng Hang]{Cheng Hang} % Your name
\institute[DUFE] % Your institution as it will appear on the bottom of every slide, may be shorthand to save space
{School of Fintech \\
Dongbei University of Finance \& Economics \\ % Your institution for the title page
\medskip
\textbf{chenghangch@qq.com}% Your email address
}
\date{\today} % Date, can be changed to a custom date

	\AtBeginSection[]
{
	\begin{frame}{Content}
		\tableofcontents[sectionstyle=show/shaded,subsectionstyle=show/shaded/hide] %这几个参数我也不知道该如何准确地解释，反正它们最终的效果是突出显示当前章节，而其它章节都进行了淡化处理
		\addtocounter{framenumber}{-1}  %目录页不计算页码
	\end{frame}
}

%----------------------------------------------------------------------------------------
%	Begin
%----------------------------------------------------------------------------------------

\begin{document}

\begin{frame}
\titlepage % Print the title page as the first slide
\end{frame}

%\begin{frame}
%\frametitle{Overview} % Table of contents slide, comment this block out to remove it
%\tableofcontents % Throughout your presentation, if you choose to use \section{} and \subsection{} commands, these will automatically be printed on this slide as an overview of your presentation
%\end{frame}

%--------------------------------------
%	PRESENTATION SLIDES
%--------------------------------------

\section{The Factor Zoo Problem}

\begin{frame}{``Factor Zoo"}
	In his 2011 Presidential address to the American Finance Association, \href{https://onlinelibrary.wiley.com/doi/10.1111/j.1540-6261.2011.01671.x}{\textcolor{blue}{John Cochrane}}  asks three questions about what he describes as the “zoo” of new anomalies:
	\begin{enumerate}
		\item Which characteristics really provide independent information about average returns?
		\item Does each new anomaly variable also correspond to a new factor formed on those same anomalies?
		\item Do many of these new factors are really important (and can account for many
		characteristics)?
	\end{enumerate}
\end{frame}

\begin{frame}{The Proliferation of Factors: A Crisis?}
    \textbf{\href{https://onlinelibrary.wiley.com/doi/10.1111/j.1540-6261.2011.01671.x}{\textcolor{blue}{Cochrane (2011)}} \textcolor{red}{Warning}:}
    \begin{itemize}
        \item ``A zoo of new factors"
        \item Threatens coherent understanding of asset pricing
        \item Calls for parsimony and economic interpretation
    \end{itemize}

    \textbf{The Explosion of Factors:}
    \begin{itemize}
        \item \textcolor{blue}{\href{https://doi.org/10.1093/rfs/hhv059}{Harvey, Liu, and Zhu (2016)}}: 316 factors documented in top journals
        \item \textcolor{blue}{\href{https://academic.oup.com/rfs/article/30/12/4389/3091648}{Green, Hand, and Zhang (2017)}}: 94 characteristics with "independent" information
        \item \textcolor{blue}{\href{https://doi.org/10.1093/rfs/hhy131}{Hou, Xue, and Zhang (2020)}}: Attempt to replicate 452 anomalies
    \end{itemize}

\end{frame}

\begin{frame}{The Statistical Crisis: Multiple Testing Problem}
    \textbf{\textcolor{blue}{\href{https://doi.org/10.1093/rfs/hhv059}{Harvey, Liu, and Zhu (2016)}} "...and the Cross-Section of Expected Returns"}
    
    
    \textbf{Key Findings:}
    \begin{itemize}
        \item 316 factors published across 313 top journal articles
        \item Traditional t-stat threshold (1.96) is too low
        \item Propose new hurdle: \textbf{t-stat $> 3.0$} for new discoveries
        \item Most recently discovered factors fail this higher standard
    \end{itemize}
    
    
    \textbf{The Multiple Testing Problem:}
    \begin{equation*}
        \text{Prob}(\text{At least one false positive}) = 1 - (1-0.05)^{316} \approx 100\%
    \end{equation*}
    
    \begin{itemize}
        \item With 316 tests at 5\% level, expect 16 false discoveries
        \item Publication bias: Only significant results get published
        \item Bonferroni correction: Divide $\alpha$ by number of tests
    \end{itemize}
    
    
    \textbf{Implications:}
    \begin{itemize}
        \item Many ``significant" factors are likely spurious
        \item Need higher standards for claiming new discoveries
        \item Out-of-sample validation crucial
    \end{itemize}
\end{frame}

\begin{frame}{The Replication Crisis}
    \textbf{\textcolor{blue}{\href{https://doi.org/10.1093/rfs/hhy131}{Hou, Xue, and Zhang (2020)}} ``Replicating Anomalies"}
    
    
    \textbf{Research Design:}
    \begin{itemize}
        \item Attempt to replicate 452 anomalies from 447 studies
        \item Use consistent methodology and data
        \item Test both equal-weighted and value-weighted portfolios
    \end{itemize}
    
    
    \textbf{Shocking Results:}
    \begin{itemize}
        \item Only \textbf{64\%} significant at 5\% level (should be 95\%)
        \item Only \textbf{42\%} significant in value-weighted portfolios
        \item \textbf{85\%} driven by microcaps (smallest 20\% of firms)
        \item Many anomalies disappear with slight methodology changes
    \end{itemize}
    
    
    \textbf{Reasons for Failure:}
    \begin{enumerate}
        \item \textbf{Microcap bias:} Many anomalies only work in illiquid stocks
        \item \textbf{Data errors:} Mistakes in original implementation
        \item \textbf{Specification choices:} Arbitrary decisions drive results
        \item \textbf{Pure data mining:} Never real to begin with
    \end{enumerate}
\end{frame}

\begin{frame}{Post-Publication Decay}
    \textbf{\textcolor{blue}{\href{https://doi.org/10.1093/rfs/hhw101}{McLean and Pontiff (2016)}} "Does Academic Research Destroy Predictability?"}
    
    
    \textbf{Key Finding:}
    \begin{itemize}
        \item Anomaly returns decline by \textbf{26\%} (equal-weighted) after publication
        \item \textbf{32\%} decline for microcaps
        \item \textbf{58\%} decline for short legs of anomalies
        \item Decay begins even before publication (working paper stage)
    \end{itemize}
    
    
    \textbf{Two Competing Explanations:}
    \begin{enumerate}
        \item \textbf{Arbitrage:} Sophisticated investors trade on published anomalies
        \begin{itemize}
            \item Partial decay suggests arbitrage is limited
            \item Short-sale constraints bind for negative alphas
        \end{itemize}
        
        \item \textbf{Data Mining:} In-sample overfitting inflates original results
        \begin{itemize}
            \item Out-of-sample returns closer to true magnitude
            \item Publication bias exaggerates effect sizes
        \end{itemize}
    \end{enumerate}
    
    
    \textbf{Evidence suggests BOTH mechanisms at work}
\end{frame}

\begin{frame}{Additional Statistical Issues}
    \textbf{1. P-Hacking (\textcolor{blue}{\href{https://doi.org/10.1111/jofi.12530}{Harvey, 2017}}):}
    \begin{itemize}
        \item Researchers try multiple specifications until finding significance
        \item File drawer problem: Negative results unpublished
        \item Solution: Pre-registration of hypotheses and methods
    \end{itemize}
    
    
    \textbf{2. Data Snooping (\textcolor{blue}{\href{https://academic.oup.com/rfs/article/31/7/2606/4977829}{Linnainmaa and Roberts, 2018}}):}
    \begin{itemize}
        \item Same datasets used repeatedly across studies
        \item Spurious patterns emerge from repeated testing
        \item Need fresh, out-of-sample data
    \end{itemize}
    
    
    \textbf{3. Errors-in-Variables (\textcolor{blue}{\href{https://doi.org/10.1016/j.jfineco.2019.02.010}{Jegadeesh et al., 2019}}):}
    \begin{itemize}
        \item Beta estimates contain measurement error
        \item Traditional two-pass regressions biased
        \item Propose instrumental variables approach
        \item Find: Most factor risk premiums insignificant after controlling for characteristics
    \end{itemize}
\end{frame}

\section{The Factor Wars: Competing Models}

\begin{frame}{Evolution of Multi-Factor Models}
    \begin{table}
        \scriptsize
        \begin{tabular}{p{3.5cm}p{6cm}}
            \toprule
            \textbf{Model} & \textbf{Factors} \\
            \midrule
            \textcolor{blue}{\href{https://doi.org/10.1016/0304-405X(93)90023-5}{Fama-French (1993)}} & Market, Size (SMB), Value (HML) \\[0.2cm]
            \textcolor{blue}{\href{https://doi.org/10.1111/j.1540-6261.1997.tb03808.x}{Carhart (1997)}} & FF3 + Momentum (UMD) \\[0.2cm]
            \textcolor{blue}{\href{https://doi.org/10.1016/j.jfineco.2013.01.003}{Novy-Marx (2013)}} & Market, Size, Value, Profitability \\[0.2cm]
            \textcolor{blue}{\href{https://doi.org/10.1016/j.jfineco.2014.02.010}{Fama-French (2015)}} & Market, Size, Value, Profitability (RMW), Investment (CMA) \\[0.2cm]
            \textcolor{blue}{\href{https://doi.org/10.1093/rfs/hhu068}{Hou, Xue, Zhang (2015)}} & Market, Size, Investment (I/A), Profitability (ROE) \\[0.2cm]
            \textcolor{blue}{\href{https://doi.org/10.1093/rfs/hhw107}{Stambaugh-Yuan (2017)}} & Market, Size, Management (MGMT), Performance (PERF) \\[0.2cm]
            \textcolor{blue}{\href{https://doi.org/10.1016/j.jfineco.2018.02.012}{Fama-French (2018)}} & FF5 + Momentum (FF6) \\[0.2cm]
            \textcolor{blue}{\href{https://doi.org/10.1093/rfs/hhz069}{Daniel et al. (2020)}} & Market, Financing, Post-earnings drift, Long/short-term contrarian \\
            \bottomrule
        \end{tabular}
    \end{table}
    
    
    \textit{From 3 factors (1993) to dozens of competing models (2020s)}
\end{frame}

\begin{frame}{Fama-French vs. Q-Factors: The Core Debate}
    \textbf{\textcolor{blue}{\href{https://doi.org/10.1093/rof/rfy032}{Hou, Mo, Xue, Zhang (2019)}} ``Which Factors?"}
    
    \textbf{Key Claims:}
    \begin{enumerate}
        \item \textbf{Spanning Tests:} Q-factor model subsumes FF5 and FF6
        \begin{itemize}
            \item FF factors have insignificant alphas vs. Q-factors
            \item But Q-factors have significant alphas vs. FF
        \end{itemize}
        
        \item \textbf{Construction Method Matters:}
        \begin{itemize}
            \item Stambaugh-Yuan MGMT and PERF are correlated 0.8+ with Q-factors
            \item Once replicated with traditional methods, very similar
            \item Differences driven by construction details, not economics
        \end{itemize}
        
        \item \textbf{Theoretical Support:}
        \begin{itemize}
            \item Investment CAPM predicts investment-return relation
            \item Valuation theory supports profitability effect
        \end{itemize}
    \end{enumerate}
\end{frame}


\begin{frame}{Fama-French Response: Choosing factors \href{https://doi.org/10.1016/j.jfineco.2018.02.012}{\textcolor{blue}{(Fama and French, 2018)}}}

    \begin{itemize}
        \item Core Idea: The optimal model should maximize the squared Sharpe ratio ($Sh^{2}(f)$) of its factors.
        \item Formula:
        \begin{align*}
        Sh^{2}(\alpha) = \sum_{i = 1}^{n}\frac{a_{i}^{2}}{\sigma_{i}^{2}} = a^{\prime}\Sigma^{-1}a
        \end{align*}
        where $a_{i}$ is the intercept, $\sigma_{i}$ is the residual standard deviation, and $\Sigma$ is the residual covariance matrix.
        \item Minimization Objective: Model evaluation should minimize the squared Sharpe ratio ($Sh^{2}(a)$) of the intercept vector.
        \item Model Selection: Select the model with the highest $Sh^{2}(f)$, that is, the model with the lowest $Sh^{2}(a)$.
    \end{itemize}
\end{frame}


\begin{frame}{Choosing factors: \href{https://doi.org/10.1016/j.jfineco.2018.02.012}{\textcolor{blue}{Fama and French (2018)}}}
    \begin{itemize}
\item Proposes using the \textbf{maximum squared Sharpe ratio} ($Sh^2(f)$) of model factors as a metric for ranking asset pricing models
\item \textbf{Minimize unexplained average returns} (intercepts) across all assets
\item Examines:
\begin{itemize}
\item Nested models: CAPM, 3-factor, 5-factor, 6-factor
\item Non-nested variations addressing:
\begin{enumerate}
\item Cash vs. operating profitability
\item Spread factors vs. excess returns
\item Small stock vs. combined size factors
\end{enumerate}
\end{itemize}
\end{itemize}
\end{frame}

\begin{frame}{Choosing factors: \href{https://doi.org/10.1016/j.jfineco.2018.02.012}{\textcolor{blue}{Fama and French (2018)}}}
\begin{itemize}
\item \textbf{Nested models:} 6-factor model (adding momentum) performs best
\item \textbf{Profitability:} Cash profitability ($RMW_c$) $\gg$ operating profitability ($RMW_o$) in $Sh^2 (f)$
\item \textbf{Factor construction:}
\begin{itemize}
\item Small stock factors outperform combined size factors
\item Best model: Mkt + SMB + small stock HML, RMWc, CMA, UMD
\end{itemize}
\end{itemize}
\end{frame}


\section{Model Selection: The Maximum Sharpe Ratio Criterion}

\begin{frame}{The Model Selection Problem}
    \textbf{The Challenge:}
    \begin{itemize}
        \item 316+ factors documented (Harvey et al., 2016)
        \item Multiple competing models: FF5, FF6, Q-factors, behavioral factors
        \item How to choose the "best" model?
    \end{itemize}
    
    \vspace{0.3cm}
    \textbf{Traditional Approaches Have Problems:}
    \begin{enumerate}
        \item \textbf{GRS test:} Depends on choice of test assets
        \item \textbf{Spanning tests:} Only work for nested models, often circular
        \item \textbf{Information criteria (AIC/BIC):} Arbitrary penalties
    \end{enumerate}
    
    \vspace{0.3cm}
    \textbf{The Barillas-Shanken (2017, 2018) Solution:}
    \begin{itemize}
        \item Use \textbf{maximum squared Sharpe ratio} of factors
        \item Theoretically grounded in SDF framework
        \item Works for nested and non-nested models
        \item Directly interpretable: Best risk-return tradeoff
    \end{itemize}
    
    
    \textit{Following: \textcolor{blue}{\href{https://doi.org/10.1093/rfs/hhx182}{Barillas and Shanken (2017)}}, \textcolor{blue}{\href{https://doi.org/10.1093/rfs/hhz089}{Barillas and Shanken (2018)}}}
\end{frame}

\section{Theoretical Foundation}

\begin{frame}{Starting Point: The SDF Framework}
    \textbf{Fundamental Pricing Equation:}
    \begin{equation}\label{eq:sdf_basic}
        E[m_{t+1} R_{i,t+1}] = 1 \quad \text{for all assets } i
    \end{equation}
    where $m_{t+1}$ is the stochastic discount factor (SDF)
    
    \vspace{0.3cm}
    \textbf{For Excess Returns:}
    \begin{equation}\label{eq:sdf_excess}
        E[m_{t+1} R^e_{i,t+1}] = 0 \quad \text{for all } i
    \end{equation}
    
    \vspace{0.3cm}
    \textbf{Key Properties:}
    \begin{itemize}
        \item SDF $m$ is unique if markets complete
        \item If incomplete, many valid SDFs exist
        \item All valid SDFs lie in same \textit{SDF space}
        \item Our goal: Find SDF that best prices assets
    \end{itemize}
    
    
    \textbf{Connection to Factor Models:}
    \begin{itemize}
        \item Linear factor model assumes $m = a + b' f$
        \item Different factor sets $\rightarrow$ Different SDFs
        \item Which factor set gives the "best" SDF?
    \end{itemize}
\end{frame}

\begin{frame}{The Hansen-Jagannathan Distance}
    \textbf{Measuring SDF Quality:}
    
    For a candidate SDF $m^*$, define the pricing error:
    \begin{equation*}
        e_i = E[m^* R^e_i] \neq 0 \quad \text{(if model wrong)}
    \end{equation*}
    
    
    \textbf{Hansen-Jagannathan (1997) Distance:}
    \begin{equation}\label{eq:hj_distance}
        HJ = \min_m \sqrt{E[(m - m^*)^2]} \quad \text{subject to } E[m R^e_i] = 0 \text{ for all } i
    \end{equation}
    
    \begin{itemize}
        \item Minimum distance from $m^*$ to set of valid SDFs
        \item If $HJ = 0$, then $m^*$ prices all assets correctly
        \item If $HJ > 0$, then $m^*$ has pricing errors
    \end{itemize}
    
    \vspace{0.3cm}
    \textbf{Key Result (Hansen-Jagannathan, 1997):}
    \begin{equation}\label{eq:hj_alpha}
        HJ^2 = e' \Sigma^{-1} e = \alpha' \Sigma_\varepsilon^{-1} \alpha
    \end{equation}
    where $\alpha$ is vector of pricing errors (intercepts), $\Sigma_\varepsilon$ is residual covariance matrix
    
    
    \textit{Minimizing HJ distance $\Leftrightarrow$ Minimizing weighted sum of squared alphas}
\end{frame}

\begin{frame}{The Hansen-Jagannathan Bound}
    \textbf{The Fundamental Inequality:}
    
    For any SDF $m$ and any return $R^e$:
    \begin{equation}\label{eq:hj_bound}
        \frac{\sigma(m)}{E[m]} \geq \frac{|E[R^e]|}{\sigma(R^e)} = |SR(R^e)|
    \end{equation}
    
    where $SR(R^e)$ is the Sharpe ratio of $R^e$
    
    \vspace{0.3cm}
    \textbf{Proof Sketch:}
    \begin{equation*}
        \begin{aligned}
            E[m R^e] &= 0 \quad \text{(fundamental equation)} \\
            \Rightarrow \text{Cov}(m, R^e) &= -E[m] E[R^e] \\
            \Rightarrow |\text{Cov}(m, R^e)| &= E[m] |E[R^e]| \\
            \Rightarrow \sigma(m) \sigma(R^e) &\geq E[m] |E[R^e]| \quad \text{(Cauchy-Schwarz)} \\
            \Rightarrow \frac{\sigma(m)}{E[m]} &\geq \frac{|E[R^e]|}{\sigma(R^e)}
        \end{aligned}
    \end{equation*}
    
    
    \textbf{Economic Interpretation:}
    \begin{itemize}
        \item LHS: Volatility of SDF relative to its mean
        \item RHS: Maximum Sharpe ratio achievable by any portfolio
        \item Bound is tight: Equality holds for tangency portfolio
    \end{itemize}
\end{frame}

\begin{frame}{The Maximum Sharpe Ratio}
    \textbf{Definition:}
    
    Consider all possible portfolios formed from assets. The \textbf{maximum squared Sharpe ratio} is:
    \begin{equation}\label{eq:max_sr}
        SR^2_{\max} = \max_w \frac{(w' \mu)^2}{w' \Sigma w}
    \end{equation}
    where $w$ is portfolio weights, $\mu = E[R^e]$ is expected excess returns, $\Sigma = \text{Var}(R^e)$
    
    \vspace{0.3cm}
    \textbf{Solution (Mean-Variance Optimization):}
    \begin{equation*}
        SR^2_{\max} = \mu' \Sigma^{-1} \mu
    \end{equation*}
    
    
    \textbf{Connection to HJ Bound:}
    \begin{itemize}
        \item From HJ bound: $\frac{\sigma(m)}{E[m]} \geq SR_{\max}$
        \item Equality holds when $m$ is in the mean-variance frontier
        \item At equality: $\left(\frac{\sigma(m)}{E[m]}\right)^2 = SR^2_{\max} = \mu' \Sigma^{-1} \mu$
    \end{itemize}
    
    
    \textbf{Key Insight:}
    \begin{itemize}
        \item $SR^2_{\max}$ measures the "sharpness" of investment opportunities
        \item Higher $SR^2_{\max}$ $\rightarrow$ Better risk-return tradeoffs available
    \end{itemize}
\end{frame}

\begin{frame}{Linear Factor Models and the SDF}
    \textbf{Factor Model Specification:}
    \begin{equation}\label{eq:factor_model}
        R^e_{i,t+1} = \alpha_i + \beta_i' f_{t+1} + \varepsilon_{i,t+1}
    \end{equation}
    
    \textbf{Implied SDF (Linear):}
    \begin{equation}\label{eq:sdf_linear}
        m_{t+1} = a + b' f_{t+1}
    \end{equation}
    
    
    \textbf{Pricing Restriction:}
    \begin{equation*}
        E[m_{t+1} R^e_{i,t+1}] = 0 \quad \Rightarrow \quad E[R^e_i] = -\frac{b}{a} E[f] = \beta_i' \lambda
    \end{equation*}
    where $\lambda = E[f]$ is the vector of factor risk premiums
    
    \vspace{0.3cm}
    \textbf{Key Properties:}
    \begin{enumerate}
        \item If $\alpha_i = 0$ for all $i$, model prices all assets correctly
        \item SDF is determined by factor premiums: $m \propto 1 - \lambda' \Sigma_f^{-1} (f - E[f])$
        \item Different factor sets $\rightarrow$ Different SDFs with different pricing accuracy
    \end{enumerate}
\end{frame}

\begin{frame}{The Equivalence Theorem}
    \textbf{Barillas-Shanken (2017) Key Result:}
    
    For \textit{traded} factors (portfolios), the following are equivalent:
    
    
    \begin{theorem}
    Let $f_A$ and $f_B$ be two sets of traded factors. Then:
    \begin{equation}\label{eq:equivalence}
        SR^2(f_A) > SR^2(f_B) \quad \Leftrightarrow \quad SR^2(\alpha_B | f_A) > SR^2(\alpha_A | f_B)
    \end{equation}
    \end{theorem}
    
    
    \textbf{In Words:}
    \begin{itemize}
        \item Model A has higher factor Sharpe ratio than B
        \item $\Leftrightarrow$ Model B's factors have larger alphas vs. A than vice versa
        \item $\Leftrightarrow$ Model A explains B better than B explains A
    \end{itemize}
    
    \vspace{0.3cm}
    \textbf{Implication:}
    \begin{itemize}
        \item Maximizing $SR^2(f)$ $\Leftrightarrow$ Minimizing $SR^2(\alpha)$ on test assets
        \item Can rank models by either metric
        \item For model selection: Choose model with highest $SR^2(f)$
    \end{itemize}
\end{frame}

\begin{frame}{Proof of Equivalence (Part 1): Setup}
    \textbf{Consider two factor models with traded factors:}
    \begin{align*}
        R^e_i &= \alpha_i^A + \beta_i^{A'} f_A + \varepsilon_i^A \\
        R^e_i &= \alpha_i^B + \beta_i^{B'} f_B + \varepsilon_i^B
    \end{align*}
    
    
    \textbf{Squared Sharpe Ratios:}
    \begin{align*}
        SR^2(f_A) &= E[f_A]' \text{Var}(f_A)^{-1} E[f_A] = \mu_A' \Sigma_A^{-1} \mu_A \\
        SR^2(f_B) &= E[f_B]' \text{Var}(f_B)^{-1} E[f_B] = \mu_B' \Sigma_B^{-1} \mu_B
    \end{align*}
    
    
    \textbf{Pricing Errors on Test Assets:}
    \begin{align*}
        SR^2(\alpha^A) &= \alpha^{A'} \Sigma_{\varepsilon^A}^{-1} \alpha^A \\
        SR^2(\alpha^B) &= \alpha^{B'} \Sigma_{\varepsilon^B}^{-1} \alpha^B
    \end{align*}
    
    
    \textbf{Goal:} Show that $SR^2(f_A) > SR^2(f_B) \Leftrightarrow SR^2(\alpha^A) < SR^2(\alpha^B)$
\end{frame}

\begin{frame}{Proof of Equivalence (Part 2): Key Step}
    \textbf{Spanning Regression:}
    
    Since factors are traded, we can regress one set on the other:
    \begin{equation}\label{eq:spanning}
        f_B = \alpha_{BA} + \Gamma f_A + u
    \end{equation}
    
    where $\alpha_{BA}$ are alphas of model B's factors vs. model A
    
    \vspace{0.3cm}
    \textbf{Decomposition of Sharpe Ratios:}
    
    By orthogonality of $f_A$ and $u$:
    \begin{equation}\label{eq:decomposition}
        SR^2(f_B) = SR^2(f_A) + SR^2(\alpha_{BA} | f_A)
    \end{equation}
    
    
    \textbf{Interpretation:}
    \begin{itemize}
        \item $SR^2(f_B)$: Total Sharpe ratio of model B factors
        \item $SR^2(f_A)$: Component explained by model A factors
        \item $SR^2(\alpha_{BA} | f_A)$: \textit{Incremental} Sharpe ratio beyond A
    \end{itemize}
    
    
    \textit{This decomposition is the key to the equivalence}
\end{frame}

\begin{frame}{Proof of Equivalence (Part 3): Conclusion}
    \textbf{From the Decomposition:}
    \begin{equation*}
        SR^2(f_B) - SR^2(f_A) = SR^2(\alpha_{BA} | f_A)
    \end{equation*}
    
    Similarly:
    \begin{equation*}
        SR^2(f_A) - SR^2(f_B) = SR^2(\alpha_{AB} | f_B)
    \end{equation*}
    
    \vspace{0.3cm}
    \textbf{Therefore:}
    \begin{equation*}
        SR^2(f_A) > SR^2(f_B) \quad \Leftrightarrow \quad SR^2(\alpha_{BA} | f_A) > 0
    \end{equation*}
    
    
    \textbf{Economic Interpretation:}
    \begin{itemize}
        \item Model A dominates B if and only if:
        \begin{itemize}
            \item B's factors have significant alphas vs. A, OR
            \item A's factors subsume B's factors
        \end{itemize}
        \item Can test by either:
        \begin{itemize}
            \item Computing factor Sharpe ratios, OR
            \item Running spanning regressions
        \end{itemize}
    \end{itemize}
    
    
    \textit{This proves the equivalence for traded factors}
\end{frame}

\begin{frame}{Extension to Pricing Test Assets}
    \textbf{From Factors to Alphas:}
    
    \textbf{Key Insight (Gibbons-Ross-Shanken, 1989):}
    
    For test assets $R^e$ and factors $f$:
    \begin{equation}\label{eq:grs}
        SR^2_{\max}(R^e, f) = SR^2(f) + SR^2(\alpha)
    \end{equation}
    
    where:
    \begin{itemize}
        \item $SR^2_{\max}(R^e, f)$: Maximum Sharpe ratio from combining test assets and factors
        \item $SR^2(f)$: Sharpe ratio of factors alone
        \item $SR^2(\alpha) = \alpha' \Sigma_\varepsilon^{-1} \alpha$: Sharpe ratio of pricing errors
    \end{itemize}
    
    \vspace{0.3cm}
    \textbf{Rearranging:}
    \begin{equation}\label{eq:alpha_sr}
        SR^2(\alpha) = SR^2_{\max}(R^e, f) - SR^2(f)
    \end{equation}
    
    
    \textbf{Interpretation:}
    \begin{itemize}
        \item $SR^2(\alpha)$ measures how much test assets add to Sharpe ratio
        \item If factors span test assets: $\alpha = 0 \Rightarrow SR^2(\alpha) = 0$
        \item If factors don't span: $SR^2(\alpha) > 0$ $\rightarrow$ Pricing errors
    \end{itemize}
\end{frame}

\begin{frame}[allowframebreaks]{The Model Selection Criterion}
    \textbf{Putting It All Together:}
    
    Given test assets $R^e$ with maximum achievable Sharpe ratio $SR^2_{\max}$:
    \begin{equation*}
        SR^2(\alpha) = SR^2_{\max} - SR^2(f)
    \end{equation*}
    
    \vspace{0.3cm}
    \textbf{Decision Rule:}
    \begin{center}
        \fbox{\parbox{0.85\textwidth}{
        \textbf{Choose the model that maximizes $SR^2(f)$}
        
        \vspace{0.1cm}
        Equivalently: Minimize $SR^2(\alpha)$
        
        \vspace{0.1cm}
        This gives the smallest pricing errors on test assets
        }}
    \end{center}
    
    \vspace{0.3cm}
    \textbf{Why This Works:}
    \begin{enumerate}
        \item Maximizing $SR^2(f)$ gets closest to $SR^2_{\max}$
        \item Leaves least "room" for unexplained returns (alphas)
        \item Minimizes Hansen-Jagannathan distance
        \item Finds SDF that best prices assets
    \end{enumerate}
    
    
    \textbf{Advantages:}
    \begin{itemize}
        \item Works for nested and non-nested models
        \item Clear economic interpretation
        \item Straightforward to compute
    \end{itemize}
\end{frame}

\begin{frame}{The Mathematics: Formal Statement}
    \textbf{Theorem (Barillas-Shanken, 2018):}
    
    
    Let $\mathcal{F}_A$ and $\mathcal{F}_B$ be two candidate factor models with traded factors. Define:
    
    \begin{itemize}
        \item $\theta_A^2 = SR^2(f_A) = \mu_A' \Sigma_A^{-1} \mu_A$
        \item $\theta_B^2 = SR^2(f_B) = \mu_B' \Sigma_B^{-1} \mu_B$
        \item $\psi_A^2 = SR^2(\alpha^A) = \alpha^{A'} \Sigma_{\varepsilon^A}^{-1} \alpha^A$
        \item $\psi_B^2 = SR^2(\alpha^B) = \alpha^{B'} \Sigma_{\varepsilon^B}^{-1} \alpha^B$
    \end{itemize}
    
    
    Then:
    \begin{equation}\label{eq:theorem}
        \boxed{
        \theta_A^2 > \theta_B^2 \quad \Leftrightarrow \quad \psi_A^2 < \psi_B^2
        }
    \end{equation}
    
    
    \textbf{Moreover:}
    \begin{equation*}
        \theta_A^2 - \theta_B^2 = \psi_B^2 - \psi_A^2
    \end{equation*}
    
    
    \textit{The gain in factor Sharpe ratio exactly equals the reduction in pricing error Sharpe ratio}
\end{frame}

\begin{frame}{Numerical Example}
    \textbf{Suppose we compare two models:}
    
    
    \textbf{Model A (FF5):} $f_A = \{\text{Mkt}, \text{SMB}, \text{HML}, \text{RMW}, \text{CMA}\}$
    \begin{itemize}
        \item $SR^2(f_A) = 0.45$
    \end{itemize}
    
    \textbf{Model B (FF6):} $f_B = \{\text{Mkt}, \text{SMB}, \text{HML}, \text{RMW}, \text{CMA}, \text{UMD}\}$
    \begin{itemize}
        \item $SR^2(f_B) = 0.58$
    \end{itemize}
    
    \vspace{0.3cm}
    \textbf{Test on 25 size-value portfolios:}
    \begin{itemize}
        \item Model A pricing errors: $SR^2(\alpha^A) = 0.25$
        \item Model B pricing errors: $SR^2(\alpha^B) = 0.12$
    \end{itemize}
    
    \vspace{0.3cm}
    \textbf{Verification of Equivalence:}
    \begin{align*}
        SR^2(f_B) - SR^2(f_A) &= 0.58 - 0.45 = 0.13 \\
        SR^2(\alpha^A) - SR^2(\alpha^B) &= 0.25 - 0.12 = 0.13 \quad \checkmark
    \end{align*}
    
    
    \textbf{Conclusion:} Model B (FF6) is superior $\rightarrow$ Adding momentum improves model
\end{frame}

\begin{frame}{Why Squared Sharpe Ratio? (Not Just Sharpe Ratio)}
    \textbf{Mathematical Reasons:}
    
    
    \begin{enumerate}
        \item \textbf{Additive Decomposition:}
        \begin{equation*}
            SR^2(f_A, f_B) = SR^2(f_A) + SR^2(\alpha_B | f_A)
        \end{equation*}
        Only works for \textit{squared} Sharpe ratios, not levels
        
        \item \textbf{Quadratic Forms:}
        \begin{equation*}
            SR^2 = \mu' \Sigma^{-1} \mu
        \end{equation*}
        Natural quadratic form appears in mean-variance optimization
        
        \item \textbf{Chi-squared Distributions:}
        \begin{equation*}
            T \cdot SR^2 \sim \chi^2_k \quad \text{asymptotically}
        \end{equation*}
        Enables straightforward hypothesis testing
    \end{enumerate}
    
    \vspace{0.3cm}
    \textbf{Economic Interpretation:}
    \begin{itemize}
        \item $SR^2$ measures \textit{curvature} of mean-variance frontier
        \item Higher curvature $\rightarrow$ Better investment opportunities
        \item Related to maximum expected utility gain
    \end{itemize}
\end{frame}

\begin{frame}{Connection to GRS Test}
    \textbf{Gibbons-Ross-Shanken (1989) Test:}
    
    Test $H_0: \alpha = 0$ using:
    \begin{equation}\label{eq:grs_test}
        GRS = \frac{T-N-K}{N} \cdot \frac{1}{1 + SR^2(f)} \cdot \alpha' \Sigma_\varepsilon^{-1} \alpha \sim F_{N, T-N-K}
    \end{equation}
    
    where $T$ = time periods, $N$ = test assets, $K$ = factors
    
    \vspace{0.3cm}
    \textbf{Relationship to $SR^2(\alpha)$:}
    \begin{equation*}
        GRS \propto SR^2(\alpha) = \alpha' \Sigma_\varepsilon^{-1} \alpha
    \end{equation*}
    
    
    \textbf{Key Insights:}
    \begin{enumerate}
        \item GRS test statistic directly related to $SR^2(\alpha)$
        \item Minimizing $SR^2(\alpha)$ $\Rightarrow$ Lower GRS statistic
        \item Maximum $SR^2(f)$ criterion consistent with GRS testing
        \item But: $SR^2$ criterion doesn't depend on choice of $N$ test assets
    \end{enumerate}
    
    
    \textit{The Sharpe ratio criterion generalizes GRS to model comparison}
\end{frame}

\begin{frame}[allowframebreaks]{Bayesian Interpretation}
    \textbf{Barillas-Shanken (2018) Bayesian Extension:}

    \textbf{Model Posterior Probability:}
    \begin{equation}\label{eq:bayesian}
        P(\mathcal{M}_A | \text{data}) \propto P(\text{data} | \mathcal{M}_A) \cdot P(\mathcal{M}_A)
    \end{equation}
    
    \textbf{Key Result:}
    
    Under diffuse priors, the Bayes factor comparing models A and B is:
    \begin{equation*}
        \frac{P(\text{data} | \mathcal{M}_A)}{P(\text{data} | \mathcal{M}_B)} \approx \left(\frac{1 + SR^2(f_B)}{1 + SR^2(f_A)}\right)^{T/2} \cdot \text{(complexity penalty)}
    \end{equation*}
    
    \vspace{0.3cm}
    \textbf{Interpretation:}
    \begin{itemize}
        \item Higher $SR^2(f)$ increases model likelihood
        \item But: Penalized by number of factors (parsimony)
        \item Trade-off between fit and complexity
        \item Simpler models preferred if $SR^2$ similar
    \end{itemize}
    
    \textbf{Practical Implication:}
    \begin{itemize}
        \item Don't just maximize $SR^2(f)$ blindly
        \item Consider model complexity
        \item Substantial $SR^2$ improvement needed to justify additional factors
    \end{itemize}
\end{frame}

\section{Fama-French (2018) Implementation}

\begin{frame}{FF (2018): Applying the Sharpe Ratio Criterion}
    \textbf{Research Design:}
    
    
    \textbf{1. Models Compared:}
    \begin{itemize}
        \item \textbf{Nested:} CAPM, FF3, FF5, FF6 (adding momentum)
        \item \textbf{Non-nested variations:}
        \begin{itemize}
            \item Operating vs. cash profitability ($\text{RMW}_o$ vs. $\text{RMW}_c$)
            \item Spread factors vs. excess returns
            \item Small-stock vs. combined size factors
        \end{itemize}
    \end{itemize}
    
    
    \textbf{2. Compute $SR^2(f)$ for Each Model:}
    \begin{equation*}
        SR^2(f) = \hat{\mu}_f' \hat{\Sigma}_f^{-1} \hat{\mu}_f
    \end{equation*}
    using sample estimates $\hat{\mu}_f = \bar{f}$, $\hat{\Sigma}_f = \text{sample covariance}$
    
    
    \textbf{3. Rank Models by $SR^2(f)$:}
    \begin{itemize}
        \item Higher $SR^2$ $\rightarrow$ Better model
        \item Check if differences statistically significant
    \end{itemize}
\end{frame}

\begin{frame}{FF (2018): Main Results}
    \textbf{Nested Model Comparison (Table 1 in paper):}
    
    
    \begin{table}
        \small
        \begin{tabular}{lcc}
            \toprule
            \textbf{Model} & \textbf{$SR^2(f)$} & \textbf{$\Delta SR^2$} \\
            \midrule
            CAPM & 0.14 & --- \\
            FF3 & 0.32 & +0.18 \\
            FF5 & 0.39 & +0.07 \\
            FF6 & 0.52 & +0.13 \\
            \bottomrule
        \end{tabular}
    \end{table}
    
    \vspace{0.3cm}
    \textbf{Interpretation:}
    \begin{itemize}
        \item Each nested extension improves $SR^2(f)$
        \item Momentum (UMD) adds substantial value: $\Delta SR^2 = 0.13$
        \item FF6 is best among nested models
    \end{itemize}
    
    
    \textbf{Statistical Significance:}
    \begin{itemize}
        \item Use spanning regressions to test differences
        \item Momentum has significant alpha vs. FF5: $t > 3$
        \item Conclusion: Momentum is not subsumed by FF5
    \end{itemize}
\end{frame}

\begin{frame}{FF (2018): Non-Nested Comparisons}
    \textbf{Cash vs. Operating Profitability:}
    
    
    \begin{table}
        \small
        \begin{tabular}{lc}
            \toprule
            \textbf{Model} & \textbf{$SR^2(f)$} \\
            \midrule
            FF5 with $\text{RMW}_o$ (operating) & 0.39 \\
            FF5 with $\text{RMW}_c$ (cash) & 0.47 \\
            \midrule
            Difference & +0.08 \\
            \bottomrule
        \end{tabular}
    \end{table}
    
    
    \textbf{Finding:} Cash profitability substantially outperforms operating profitability
    
    \vspace{0.3cm}
    \textbf{Small-Stock vs. Combined Factors:}
    \begin{itemize}
        \item \textbf{Combined:} Single HML, RMW, CMA across all size groups
        \item \textbf{Small-stock:} Separate factors for small firms only
        \item Result: Small-stock factors have higher $SR^2$ by 0.05
    \end{itemize}
    
    
    \textbf{Best Overall Model:}
    \begin{itemize}
        \item Market + SMB + Small-stock $\text{HML}_c$, $\text{RMW}_c$, CMA + UMD
        \item $SR^2(f) = 0.55$
    \end{itemize}
\end{frame}

\begin{frame}{Practical Implementation: Step-by-Step}
    \textbf{How to Apply the Sharpe Ratio Criterion:}
    
    
    \textbf{Step 1: Compute Sample Moments}
    \begin{itemize}
        \item Mean factor returns: $\hat{\mu}_f = \frac{1}{T} \sum_{t=1}^T f_t$
        \item Factor covariance: $\hat{\Sigma}_f = \frac{1}{T-1} \sum_{t=1}^T (f_t - \hat{\mu}_f)(f_t - \hat{\mu}_f)'$
    \end{itemize}
    
    
    \textbf{Step 2: Compute $SR^2(f)$}
    \begin{equation*}
        \widehat{SR^2}(f) = \hat{\mu}_f' \hat{\Sigma}_f^{-1} \hat{\mu}_f
    \end{equation*}
    
    
    \textbf{Step 3: Compare Across Models}
    \begin{itemize}
        \item Rank by $\widehat{SR^2}(f)$
        \item Test significance via spanning regressions
    \end{itemize}
    
    
    \textbf{Step 4: Verify with Test Assets (Optional)}
    \begin{itemize}
        \item Regress test portfolios on factors
        \item Compute $\widehat{SR^2}(\alpha) = \hat{\alpha}' \hat{\Sigma}_\varepsilon^{-1} \hat{\alpha}$
        \item Verify: Higher $SR^2(f)$ $\Leftrightarrow$ Lower $SR^2(\alpha)$
    \end{itemize}
\end{frame}

\begin{frame}{Advantages and Limitations}
    \textbf{Advantages of $SR^2$ Criterion:}
    
    
    \begin{enumerate}
        \item \textbf{Theoretically grounded:} Based on SDF/mean-variance theory
        \item \textbf{General:} Works for nested and non-nested models
        \item \textbf{Interpretable:} Measures risk-return efficiency
        \item \textbf{Asset-independent:} Doesn't depend on choice of test portfolios
        \item \textbf{Consistent:} Equivalent to minimizing pricing errors
    \end{enumerate}
    
    \vspace{0.3cm}
    \textbf{Limitations:}
    
    
    \begin{enumerate}
        \item \textbf{Requires traded factors:} Non-traded factors need different treatment
        \item \textbf{Sample uncertainty:} $\widehat{SR^2}$ has large standard errors
        \item \textbf{No complexity penalty:} Unlike Bayesian approach
        \item \textbf{Mean-variance framework:} Assumes quadratic utility/normal returns
        \item \textbf{Ignores factor interpretability:} Focuses only on statistical fit
    \end{enumerate}
    
    
    \textbf{Best Practice:}
    \begin{itemize}
        \item Use $SR^2$ as primary criterion
        \item But also consider: Economic interpretation, parsimony, robustness
    \end{itemize}
\end{frame}

\begin{frame}{Summary: Why Maximum Sharpe Ratio Works}
    \textbf{The Logic Chain:}
    
    
    \begin{enumerate}
        \item Assets must be priced by some SDF: $E[m R^e] = 0$
        
        \item Hansen-Jagannathan bound: $\frac{\sigma(m)}{E[m]} \geq SR_{\max}$
        
        \item Equality holds for mean-variance efficient SDF
        
        \item Linear factor models imply specific SDF: $m = a + b' f$
        
        \item Better factors $\rightarrow$ Higher $SR^2(f)$ $\rightarrow$ SDF closer to efficient frontier
        
        \item By GRS decomposition: $SR^2_{\max} = SR^2(f) + SR^2(\alpha)$
        
        \item Therefore: $\max SR^2(f) \Leftrightarrow \min SR^2(\alpha)$
        
        \item Conclusion: Choose model with highest factor Sharpe ratio
    \end{enumerate}
    
    \vspace{0.3cm}
    \textbf{Bottom Line:}
    \begin{itemize}
        \item Maximizing $SR^2(f)$ finds factors that best approximate true SDF
        \item Minimizes unexplained returns (pricing errors)
        \item Provides clear, economically interpretable ranking of models
    \end{itemize}
\end{frame}

\begin{frame}{The Verdict: No Clear Winner}
    \textbf{Cross-Model Comparisons:}
    
    \begin{table}
        \footnotesize
        \begin{tabular}{p{3cm}|p{2.5cm}|p{2.5cm}}
            \hline
            \textbf{Criterion} & \textbf{Winner} & \textbf{Source} \\
            \hline
            Squared Sharpe ratio & FF6 & Fama-French (2018) \\
            Spanning tests & Q-factors & Hou et al. (2019) \\
            Anomaly alphas & Mixed & Multiple studies \\
            Bayesian posterior & Simpler models & Barillas-Shanken (2018) \\
            \hline
        \end{tabular}
    \end{table}
    
    
    \textbf{Why No Consensus?}
    \begin{enumerate}
        \item \textbf{Test asset dependence:} Different results for different portfolios
        \item \textbf{Sample period sensitivity:} 1960s vs. recent decades
        \item \textbf{Weighting schemes:} Equal-weighted vs. value-weighted
        \item \textbf{Microcap effects:} Include or exclude smallest stocks?
        \item \textbf{Factor construction:} Breakpoints, rebalancing frequency matter
    \end{enumerate}
    
    
    \textbf{Practical Implications:}
    \begin{itemize}
        \item Report results for multiple models
        \item Be transparent about robustness (or lack thereof)
        \item Consider economic mechanism, not just $R^2$
    \end{itemize}
\end{frame}

\section{New Methodologies and Solutions}

\begin{frame}{Machine Learning Approaches}
    \textbf{The Promise: Handle High Dimensions}
    \begin{itemize}
        \item Traditional methods break down with 100+ characteristics
        \item ML techniques designed for $p \gg n$ problems
        \item Can capture non-linearities and interactions
    \end{itemize}
    
    
    \textbf{1. IPCA: Instrumented Principal Component Analysis}
    
    \textbf{\textcolor{blue}{\href{https://doi.org/10.1016/j.jfineco.2019.05.001}{Kelly, Pruitt, and Su (2019)}}}
    \begin{itemize}
        \item Allow factor loadings to vary with firm characteristics
        \item $r_{it} = \beta_{it}' f_t + \varepsilon_{it}$, where $\beta_{it} = \Gamma' z_{it}$
        \item Use PCA to extract latent factors with time-varying loadings
        \item Characteristics instrument for unobservable risk exposures
    \end{itemize}
    
    
    \textbf{Key Result:}
    \begin{itemize}
        \item Finds 5 latent factors explain most of cross-section
        \item If characteristics capture risk, IPCA should work well
        \item Outperforms fixed-loading models (FF5, Q-factors)
    \end{itemize}
\end{frame}

\begin{frame}{Machine Learning Approaches (cont.)}
    \textbf{2. LASSO and Regularization Methods}
    
    \textbf{\textcolor{blue}{\href{https://doi.org/10.1111/jofi.12883}{Feng, Giglio, and Xiu (2020)}} ``Taming the Factor Zoo"}
    \begin{itemize}
        \item Apply three-pass regression filter (3PRF) with LASSO
        \item \textbf{Pass 1:} Time-series regressions to estimate loadings
        \item \textbf{Pass 2:} Cross-sectional regressions with LASSO penalty
        \begin{equation*}
            \min_{\lambda} \sum_{t=1}^T (r_t - \beta_t \lambda)' (r_t - \beta_t \lambda) + \eta ||\lambda||_1
        \end{equation*}
        \item \textbf{Pass 3:} Re-estimate without penalty on selected factors
    \end{itemize}
    
    
    \textbf{Key Findings:}
    \begin{itemize}
        \item Of 150+ factors, LASSO selects only 10-15 as relevant
        \item Selected factors change over time (adaptive)
        \item Out-of-sample $R^2$ higher than kitchen-sink models
    \end{itemize}
    
    
    \textbf{3. Other ML Methods:}
    \begin{itemize}
        \item \textcolor{blue}{\href{https://doi.org/10.1093/rfs/hhaa009}{Gu, Kelly, and Xiu (2020)}}: Neural networks
        \item Random forests for non-linear interactions
        \item Elastic net (LASSO + Ridge)
    \end{itemize}
\end{frame}

\begin{frame}{Testing for Omitted Factors}
    \textbf{The Residual Structure Test:}
    
    \textbf{\textcolor{blue}{\href{https://doi.org/10.1016/j.jeconom.2019.06.001}{Gagliardini, Ossola, and Scaillet (2019)}}}
    
    
    \textbf{Core Idea:}
    \begin{itemize}
        \item If model is complete, residuals should have no factor structure
        \item Test by examining eigenvalues of residual covariance matrix
        \item Large eigenvalues $\rightarrow$ Missing common factors
    \end{itemize}
    
    
    \textbf{Test Statistic:}
    \begin{equation*}
        \text{Eigenvalue ratio} = \frac{\lambda_{\max}(\hat{\Sigma}_\varepsilon)}{\text{trace}(\hat{\Sigma}_\varepsilon)/n}
    \end{equation*}
    
    \begin{itemize}
        \item Under null (no omitted factors): Ratio bounded
        \item Under alternative: Ratio diverges as $n \to \infty$
    \end{itemize}
    
    
    \textbf{Application:}
    \begin{itemize}
        \item FF5: \textbf{Reject} null, missing factors remain
        \item Q-factor: \textbf{Reject} null as well
        \item Adding more factors reduces but doesn't eliminate residual structure
    \end{itemize}
\end{frame}

\begin{frame}{SDF-Based Approaches}
    \textbf{The Fundamental Pricing Equation:}
    \begin{equation*}
        E[m_{t+1} R_{t+1}] = 1 \quad \text{for all returns } R
    \end{equation*}
    
    \textbf{Hansen-Jagannathan Bound:}
    \begin{equation*}
        \frac{\sigma(m)}{E[m]} \geq \frac{|E[R^e]|}{\sigma(R^e)} = \text{Sharpe ratio}
    \end{equation*}
    
    
    \textbf{\textcolor{blue}{\href{https://doi.org/10.1086/714090}{Giglio and Xiu (2021)}} ``Asset Pricing with Omitted Factors"}
    \begin{itemize}
        \item Estimate risk premiums when some factors unobserved
        \item Use three-pass regression with bias correction
        \item Recover SDF even with incomplete factor set
    \end{itemize}
    
    
    \textbf{\textcolor{blue}{\href{https://doi.org/10.1093/rfs/hhaa111}{Giglio, Liao, and Xiu (2020)}} ``Thousands of Alphas"}
    \begin{itemize}
        \item Test thousands of anomalies simultaneously
        \item Control for data mining using latent factor structure
        \item Find: Few anomalies survive after proper controls
    \end{itemize}
\end{frame}

\begin{frame}{Accounting for Measurement Error}
    \textbf{\textcolor{blue}{\href{https://doi.org/10.1016/j.jfineco.2019.02.010}{Jegadeesh, Noh, Pukthuanthong, Roll, Wang (2019)}}}
    
    
    \textbf{The EIV Problem in Asset Pricing:}
    \begin{itemize}
        \item Factor betas estimated with error: $\hat{\beta} = \beta + \eta$
        \item Standard two-pass regression: $r_i = \lambda_0 + \lambda' \hat{\beta}_i + \varepsilon_i$
        \item EIV bias: $\hat{\lambda}$ attenuated toward zero
    \end{itemize}
    
    
    \textbf{Instrumental Variables Solution:}
    \begin{itemize}
        \item Use characteristics as instruments for betas
        \item Characteristics correlated with true betas but not measurement error
        \item Two-stage least squares (2SLS) estimation
    \end{itemize}
    
    
    \textbf{Shocking Result:}
    \begin{itemize}
        \item After correcting for EIV bias:
        \begin{itemize}
            \item CAPM market premium: \textbf{Insignificant}
            \item FF3 factor premiums: \textbf{Mostly insignificant}
            \item FF5 factor premiums: \textbf{Insignificant}
        \end{itemize}
        \item But characteristics still predict returns strongly
        \item Interpretation: Characteristics matter, not covariances (again!)
    \end{itemize}
\end{frame}

\begin{frame}{The Characteristics Response}
    \textbf{\textcolor{blue}{\href{https://doi.org/10.1093/rfs/hhaa018}{Chordia, Goyal, and Saretto (2020)}} and \textcolor{blue}{\href{https://doi.org/10.1093/rfs/hhz089}{Fama and French (2020)}}}
    
    
    \textbf{Pushback Against Jegadeesh et al. (2019):}
    \begin{enumerate}
        \item \textbf{Characteristics naturally predict better:}
        \begin{itemize}
            \item Characteristics measured precisely
            \item Betas estimated with large standard errors
            \item Of course characteristics win in horse races
        \end{itemize}
        
        \item \textbf{Doesn't prove mispricing:}
        \begin{itemize}
            \item Characteristics could still proxy for risk
            \item Through time-varying exposures
            \item Or non-linear factor loadings
        \end{itemize}
        
        \item \textbf{Economic interpretation:}
        \begin{itemize}
            \item Small firms have high betas $\rightarrow$ High risk $\rightarrow$ High returns
            \item But characteristic (size) more stable than beta estimates
            \item Use characteristic as better proxy for risk
        \end{itemize}
    \end{enumerate}
    
    
    \textbf{Bottom Line:}
    \begin{itemize}
        \item Debate unresolved
        \item Both sides make valid points
        \item May be impossible to definitively separate risk from mispricing
    \end{itemize}
\end{frame}

\section{Recent Developments (2020-2025)}

\begin{frame}{Transaction Costs and Tradability}
    \textbf{\textcolor{blue}{\href{https://doi.org/10.1093/rof/rfae034}{Johansson, Sabbatucci, and Tamoni (2025)}} ``Tradable Factor Risk Premia"}
    
    
    \textbf{The Tradability Problem:}
    \begin{itemize}
        \item Most factor portfolios have high turnover
        \item Transaction costs eat into returns
        \item Many anomalies unprofitable net of costs
    \end{itemize}
    
    
    \textbf{Key Innovations:}
    \begin{enumerate}
        \item \textbf{Construct low-turnover factor portfolios:}
        \begin{itemize}
            \item Use coarser sorts (quintiles instead of deciles)
            \item Rebalance quarterly instead of monthly
            \item Skip small buffer zones around breakpoints
        \end{itemize}
        
        \item \textbf{Estimate realistic transaction costs:}
        \begin{itemize}
            \item Bid-ask spreads from CRSP
            \item Price impact from trade sizes
            \item Different costs for institutional vs. retail
        \end{itemize}
        
        \item \textbf{Net return factors:}
        \begin{itemize}
            \item Gross return minus transaction costs
            \item More realistic for practical implementation
        \end{itemize}
    \end{enumerate}
\end{frame}

\begin{frame}{Tradable Factors: Main Results}
    \textbf{Findings (Johansson et al., 2025):}
    
    
    \begin{enumerate}
        \item \textbf{Microcap anomalies disappear:}
        \begin{itemize}
            \item High transaction costs eliminate alphas
            \item Even for institutional investors
        \end{itemize}
        
        \item \textbf{Quality and profitability survive:}
        \begin{itemize}
            \item Low turnover strategies
            \item Focus on large, liquid stocks
            \item Net returns still positive
        \end{itemize}
        
        \item \textbf{Momentum partially survives:}
        \begin{itemize}
            \item High turnover hurts
            \item But effect large enough to overcome costs
            \item Depends on implementation details
        \end{itemize}
        
        \item \textbf{Value struggles:}
        \begin{itemize}
            \item Moderate turnover + declining gross returns
            \item Net returns close to zero in recent decades
        \end{itemize}
    \end{enumerate}
    
    
    \textbf{Implication:}
    \begin{itemize}
        \item Academic factors $\neq$ Investable strategies
        \item Need to report net returns in research
    \end{itemize}
\end{frame}

\begin{frame}{Alternative Data and New Signals}
    \textbf{Beyond Traditional Accounting Data:}
    
    
    \textbf{1. \textcolor{blue}{\href{https://doi.org/10.1017/S0022109024000528}{So, Wang, and Zhang (2024)}} ``Investor Visits"}
    \begin{itemize}
        \item Track institutional investor visits to firms (unique Chinese data)
        \item Firms with more visits earn higher subsequent returns
        \item Interpretation: Visits reveal private information about firm prospects
        \item Reduces information asymmetry $\rightarrow$ Lower cost of capital
    \end{itemize}
    
    
    \textbf{2. Other Alternative Data:}
    \begin{itemize}
        \item \textbf{Satellite imagery:} Parking lot traffic, factory activity
        \item \textbf{Credit card data:} Consumer spending in real-time
        \item \textbf{Web scraping:} Product reviews, job postings
        \item \textbf{Social media:} Sentiment, attention, news flow
    \end{itemize}
    
    
    \textbf{Challenges:}
    \begin{itemize}
        \item Data mining concerns even worse with alternative data
        \item Decay after publication likely faster
        \item Expensive, proprietary access limits replication
    \end{itemize}
\end{frame}

\begin{frame}{The Current State: Factor Consolidation?}
    \textbf{Signs of Consolidation:}
    \begin{enumerate}
        \item \textbf{Core factors emerging:}
        \begin{itemize}
            \item Market, size, value, momentum, profitability, investment
            \item Most other "factors" are combinations of these
        \end{itemize}
        
        \item \textbf{Higher standards for new factors:}
        \begin{itemize}
            \item t-stat $> 3.0$ threshold widely adopted
            \item Journals requiring out-of-sample tests
            \item Pre-registration gaining traction
        \end{itemize}
        
        \item \textbf{Focus on implementation:}
        \begin{itemize}
            \item Transaction costs becoming standard
            \item Capacity constraints discussed
            \item Practitioner involvement increasing
        \end{itemize}
    \end{enumerate}
    
    
    \textbf{But Proliferation Continues:}
    \begin{itemize}
        \item Machine learning finding new patterns
        \item Alternative data sources expanding
        \item Behavioral explanations multiplying
    \end{itemize}
    
    
    \textit{The zoo is under better management, but still growing}
\end{frame}

\section{Practical Guidance and Conclusions}

\begin{frame}{Best Practices for Researchers}
    \textbf{When Proposing New Factors:}
    
    
    \textbf{1. Statistical Rigor:}
    \begin{itemize}
        \item Use t-stat $> 3.0$ threshold (Harvey et al., 2016)
        \item Report Bonferroni or other multiple testing corrections
        \item Show results are robust to:
        \begin{itemize}
            \item Different sample periods
            \item Equal vs. value-weighted portfolios
            \item Inclusion/exclusion of microcaps
            \item Various breakpoint choices
        \end{itemize}
    \end{itemize}
    
    
    \textbf{2. Out-of-Sample Validation:}
    \begin{itemize}
        \item Test in different time periods (pre-1960s if possible)
        \item International evidence (developed and emerging markets)
        \item Different asset classes (bonds, currencies, commodities)
        \item Distinguish in-sample from out-of-sample clearly
    \end{itemize}
    
    
    \textbf{3. Economic Content:}
    \begin{itemize}
        \item Provide theoretical motivation (risk or behavioral)
        \item Estimate realistic transaction costs
        \item Discuss capacity and scalability
        \item Who can actually exploit this anomaly?
    \end{itemize}
\end{frame}

\begin{frame}{Best Practices for Researchers (cont.)}
    \textbf{4. Model Testing:}
    \begin{itemize}
        \item Report alphas relative to multiple models:
        \begin{itemize}
            \item CAPM (baseline)
            \item FF3 (established benchmark)
            \item FF5 or FF6 (current standard)
            \item Q-factors (theory-based alternative)
        \end{itemize}
        \item Test whether your factor subsumes existing factors
        \item Test whether existing factors subsume yours
    \end{itemize}
    
    
    \textbf{5. Transparency and Replication:}
    \begin{itemize}
        \item Share code and data
        \item Document all data cleaning steps
        \item Explain treatment of missing data, outliers
        \item Report summary statistics of key variables
    \end{itemize}
    
    
    \textbf{6. Avoid Common Pitfalls:}
    \begin{itemize}
        \item Don't just add your factor to FF5 and declare victory
        \item Don't ignore that your factor may be combination of existing ones
        \item Don't claim "new" factor without checking related work
        \item Don't oversell implications ("death of CAPM", "solution to asset pricing")
    \end{itemize}
\end{frame}

\begin{frame}[allowframebreaks]{For Practitioners and Investors}
    \textbf{How to Evaluate Factor Strategies:}
    
    \textbf{1. Scrutinize Backtests:}
    \begin{itemize}
        \item What sample period? (In-sample overfitting?)
        \item Microcaps included? (Tradability?)
        \item Realistic transaction costs? (Implementation drag?)
        \item How selected test assets? (Cherry-picking?)
    \end{itemize}
    
    
    \textbf{2. Check Post-Publication Performance:}
    \begin{itemize}
        \item Expect 26-32\% return decay after publication
        \item Has factor "worked" since academic discovery?
        \item Arbitrage competition may have eroded returns
    \end{itemize}
    
    
    \textbf{3. Understand Economic Rationale:}
    \begin{itemize}
        \item Is this compensation for risk?
        \item Is this exploiting behavioral bias?
        \item Is this just data mining?
        \item Why should it persist going forward?
    \end{itemize}
    
    
    \textbf{4. Diversify Factor Exposure:}
    \begin{itemize}
        \item Don't bet everything on one factor
        \item Combine factors with low correlations
        \item Adjust for capacity constraints
    \end{itemize}
\end{frame}

\begin{frame}[allowframebreaks]{The Path Forward}
    \textbf{Where is Asset Pricing Headed?}
    
    
    \textbf{1. Consolidation vs. Proliferation:}
    \begin{itemize}
        \item Core factors likely to survive: market, size, value, momentum, profitability
        \item But new factors will continue to emerge
        \item Machine learning accelerating discovery
    \end{itemize}
    
    
    \textbf{2. Theory vs. Empirics:}
    \begin{itemize}
        \item Need stronger theoretical foundations
        \item But pure theory often fails empirically
        \item Middle ground: Structural models with realistic frictions
    \end{itemize}
    
    
    \textbf{3. Risk vs. Mispricing:}
    \begin{itemize}
        \item Debate unlikely to be fully resolved
        \item Both forces probably at work
        \item Context-dependent, varies by anomaly
    \end{itemize}
    
    
    \textbf{4. Academic vs. Practitioner Divide:}
    \begin{itemize}
        \item Growing convergence on implementation issues
        \item Transaction costs, capacity, market impact now standard
        \item But academic incentives still favor novelty over replication
    \end{itemize}
\end{frame}

\begin{frame}{Conclusions: What Have We Learned?}
    \textbf{Key Takeaways from the Factor Wars:}
    
    
    \begin{enumerate}
        \item \textbf{The zoo is real, but manageable:}
        \begin{itemize}
            \item Many factors are redundant or spurious
            \item Core set of 5-10 factors likely sufficient
            \item But identifying which ones remains contentious
        \end{itemize}
        
        \item \textbf{No single model dominates:}
        \begin{itemize}
            \item FF6, Q-factors, behavioral factors all have merits
            \item Performance varies by test assets, period, weighting
            \item Use multiple models, report robustness
        \end{itemize}
        
        \item \textbf{Characteristics matter more than covariances:}
        \begin{itemize}
            \item But interpretation remains ambiguous
            \item Could be risk through non-linear channels
            \item Or could be persistent mispricing
        \end{itemize}
        
        \item \textbf{Implementation details are crucial:}
        \begin{itemize}
            \item Transaction costs eliminate many anomalies
            \item Microcap bias pervasive in literature
            \item Academic backtests often not investable
        \end{itemize}
    \end{enumerate}
\end{frame}

\begin{frame}{Final Thoughts}
    \textbf{The State of Asset Pricing in 2025:}
    
    \begin{itemize}
        \item \textbf{Progress:} Better understanding of what drives returns
        \item \textbf{Humility:} Recognition of many past "discoveries" were spurious
        \item \textbf{Rigor:} Higher statistical standards being adopted
        \item \textbf{Realism:} More attention to implementation constraints
    \end{itemize}
    
    
    \textbf{Unresolved Questions:}
    \begin{itemize}
        \item Fundamental nature of alpha (risk vs. mispricing)
        \item ``Correct" number of factors
        \item Role of machine learning (help or harm?)
        \item Declining returns to factors over time
    \end{itemize}
    
    
    \textbf{For Students and Researchers:}
    \begin{itemize}
        \item Critical thinking about published results
        \item Emphasis on replication and robustness
        \item Economic intuition alongside statistical tests
        \item Bridge academic research and practical implementation
    \end{itemize}
    
    
    \textit{The factor zoo is not a crisis, but an opportunity to deepen our understanding of asset prices}
\end{frame}


\end{document} 
